\documentclass[10pt]{article}
\usepackage[margin=1in]{geometry}
\usepackage{amsmath,amssymb,booktabs,listings,xcolor,hyperref}
\hypersetup{colorlinks=true,urlcolor=blue,citecolor=blue,linkcolor=blue}
\lstset{basicstyle=\ttfamily\small,breaklines=true,columns=fullflexible,frame=single}
\title{From Teaching Examples to Repeatable Performance Experiments\\
through Comment-Directed Build Generation}
\author{Research design and prototype draft}
\date{\today}
\begin{document}
\maketitle

\begin{abstract}
Teaching examples explain an algorithm's public interface, whereas performance
drivers control setup, warmup, repetition, and reporting. Maintaining both can
introduce differences in workloads and timing boundaries. We investigate a build
workflow in which source comments identify a repeatable operation and its nested
phases. CMake retains ordinary targets and generates separate profiling sources,
libraries, and executables. A generated harness repeats the selected operation,
and Caliper records phase and repetition context. The prototype supports serial,
OpenMP, and collective MPI workloads and provides a source-linked dashboard.
This draft specifies the implementation and an evaluation protocol. The local
validation exercises functionality; it does not establish low overhead,
performance portability, or a publication-level novelty claim.
\end{abstract}

\section{Problem and research questions}
Performance experiments need an explicit workload and timing boundary that remain
consistent with the code developers read and maintain. A teaching example already
loads inputs, invokes a public algorithm, and checks its result, but its single
execution is insufficient for characterizing runtime variability. Copying it into
a performance driver creates another place to maintain those choices.

The research question is whether a small, checked annotation language and build
transformation can preserve one workload definition while providing repeatable
overall measurements and internal explanations. We evaluate four narrower questions:
\begin{enumerate}
\item Does generated execution preserve results, ownership, and supported control flow?
\item How closely do its timing boundaries match explicit hand-written drivers?
\item What runtime, build-time, and storage costs arise from the harness and phase recording?
\item Does one workflow remain useful across serial, OpenMP, and MPI workloads?
\end{enumerate}
These are hypotheses to test. Comment syntax alone is not sufficient evidence of
a new research contribution.

\section{Design and implementation}
The build transformation separates the maintained program from generated measurement
code. Python validates comments before updating any generated file. CMake custom
commands create a mirror under the build directory and regenerate it when sources,
headers, or the generator change. Original targets continue to compile the original
files. Source mappings and \texttt{\#line} directives retain navigation and compiler
diagnostics at the original locations.

\subsection{Annotation contract}
The annotations identify configuration, a repeatable value-producing operation,
and nested phases. For example, the distributed driver prepares DCSC operands before
the following block:
\begin{lstlisting}[language=C++]
const DistSpCOO<IT, NT, OT> product = [&] {
  // @profiler-benchmark time_to_solution start
  return Mpi2DSummaSpGEMM<PlusTimesRing<NT>>(a, a);
  // @profiler-benchmark time_to_solution end
}();
\end{lstlisting}
The normal executable invokes this operation once. The generated executable replaces
the annotated body with a call to a repetition harness. Its callable returns an
owning value; the harness releases the previous result outside the timed interval
and retains the last result for the example's independent reference check. Input
loading, partitioning, and conversion occur once outside this block. The selected
operation must produce a fresh result without irreversibly modifying its inputs.
The prototype does not infer or restore mutated input state.

\texttt{@profiler-region NAME start/end} identifies nested phase boundaries.
\texttt{@profiler-iter VARIABLE start/end} associates an algorithm iteration with
its region events. Benchmark repetition and algorithm iteration are distinct fields.
\texttt{@profiler-config argc argv} consumes profiling-only flags before the
application's own parser runs. Normal compilation treats all annotations as comments.

The checker tokenizes comments, literals, braces, and preprocessor branches. It
rejects unknown directives, crossed or unclosed pairs, and scope or branch crossings.
It is not a C++ parser: the compiler checks types and return paths. Region guards
close through C++ object lifetime. A benchmark block must be outside other annotated
regions; a block that returns no value is unsupported by this initial harness.

\subsection{Build products and measurement modes}
Separate library targets allow the same project to retain its ordinary build while
offering detailed instrumentation. \texttt{SPCraft} compiles original sources;
\texttt{SPCraft\_profile} compiles their generated counterparts and links Caliper.
Each selected example has three executable targets:
\begin{description}
\item[Ordinary example.] Original source and original library, executing once.
\item[Timing-only benchmark.] Generated repetition harness with the original library.
Nested annotation hooks in generated driver code are disabled at compile time.
\item[Detailed profile.] Generated harness and instrumented library, recording nested phases.
\end{description}
The timing-only executable still links Caliper to record one outer result per rank
and repetition. Its raw call timer excludes the outer annotation's begin/end and
trace output. This mode is not an instrumentation-free executable, and recording
between repetitions can affect cache and scheduling state. Comparison with a separate
hand-written reference driver remains necessary.

\subsection{Execution protocol}
The runtime controls warmups and limits without rebuilding the application.
Defaults are 1 warmup and termination after either 1000 measured repetitions or
10 seconds of batch wall time. A single explicit repetition or duration limit
replaces the other default; providing both terminates at the first limit. A running
operation completes before a duration cutoff is checked.

Warmup phase hooks and default-channel recording are disabled. For measured MPI
repetitions, the harness releases the preceding output and calls a world-communicator
barrier before starting the local call timer. The call includes algorithm-internal
allocation, communication, computation, and returned-result construction. Reductions,
stopping decisions, and outer trace recording are outside the raw call timer. Rank 0
decides whether to stop and broadcasts that decision, ensuring matching collective
execution. This prototype assumes the measured operation uses all ranks in
\texttt{MPI\_COMM\_WORLD}; subcommunicators require an explicit extension.

The outer elapsed duration is stored as integer nanoseconds. This avoids decimal
rounding of short floating-point durations during JSON export. Caliper supplies
nested phase durations and rank context. Asynchronous accelerator operations would
require explicit completion before stopping the timer and are not evaluated here.

\section{Measurement and statistical model}
Distributed timing requires reductions in a specified order. Let $t_{j,r,p,k}$ be
the inclusive duration of call $k$ to phase $p$ on rank $r$ during measured repetition
$j$. The phase sample is
\[
 T_{j,p}=\max_r\left(\sum_k t_{j,r,p,k}\right).
\]
The dashboard reports the median, mean, minimum, maximum, and sample standard
deviation of $T_{j,p}$ across repetitions. This differs from first accumulating all
repetitions per rank and then taking one maximum. Algorithm iterations can be filtered
before the per-rank sum while benchmark repetition identifiers remain intact.

Overall time to solution comes from the explicit outer call timer, with the same
per-repetition rank maximum. It is not reconstructed from child phases. Different
phases can attain their maxima on different ranks, and inclusive parent/child timers
overlap. Consequently, adding phase maxima does not produce overall elapsed time.
The dashboard retains a separate rank-total view for diagnosing imbalance and labels
those values as accumulated times.

Every expected rank must provide an outer timer for a repetition to contribute to
global statistics. The dashboard excludes partially covered repetitions and reports
their count. A genuinely unexecuted phase contributes zero when the corresponding
outer measurement is complete. The evaluation script rejects missing or duplicate
outer observations and discontinuous repetition identifiers. The final retained result
is checked against Eigen after the batch; the current prototype does not validate
every intermediate repetition.

\section{Evaluation plan and prototype validation}
The evaluation must separate functional validation from performance conclusions.
The supplied \texttt{evaluate.py} runs serial and 2-thread OpenMP SpMV and SpGEMM,
and DCSC SUMMA with 1 and 4 MPI ranks. Every case runs in timing-only and detailed
profiling modes. It verifies the numerical check, trace coverage, repetition count,
and presence or absence of nested phases. Each run retains raw traces, a portable
dashboard report, command, selected environment, input hash, binary hash, source
revision, dirty-file list, and measurement limits. Trial order is shuffled with a
recorded seed; independent launches are retained separately.

\begin{table}[ht]
\centering
\begin{tabular}{lll}
\toprule
Case & Purpose & Reference \\
\midrule
SpMV, 1 thread & Short operation and allocation boundary & Eigen \\
SpMV, 2 threads & Internally threaded workload & Eigen \\
SpGEMM, 1/2 threads & Irregular work and nested hash phases & Eigen \\
SUMMA, 1/4 ranks & Collective repetition and rank statistics & Eigen \\
Control-flow fixtures & Early exits and lexical safety & Hand-checked events \\
\bottomrule
\end{tabular}
\caption{Initial validation cases. Local execution validates the workflow; these
cases alone do not establish scaling or generality.}
\end{table}

Publication experiments should add matrices spanning shape, sparsity, and per-column
work; several problem scales; physical cores with recorded affinity; independent
process launches; and MPI execution on multiple nodes. Hand-written timing drivers
should use the identical input, output-allocation boundary, warmup policy, and
repetition schedule. A direct Caliper-API implementation should establish whether
generation changes either timing fidelity or instrumentation overhead. Google
Benchmark manual timing provides a useful shared-memory comparison with matching
boundaries~\cite{googlebenchmark}. Confidence intervals should be constructed over
independent-launch summaries rather than treating all inner repetitions as independent.

Additional experiments should measure annotation density, very short regions,
incremental rebuild cost, trace growth, missing-rank diagnostics, and injected rank
imbalance. Negative tests should cover malformed markers, invalid options, mutable
inputs, incompatible collective schedules, and exceptions. Support claims must be
limited to exercised compiler, C++ language, OpenMP, and MPI configurations.

\section{Related work and prospective contribution}
Existing tools already provide source instrumentation and rich parallel performance
analysis. Caliper supplies application annotations, contextual data, MPI-aware
measurement, and analysis services~\cite{caliper}; this prototype uses those facilities
and does not claim to introduce them. Score-P supports automatic and user-defined
instrumentation~\cite{scorepmanual}. CaPI investigates configurable selective
instrumentation~\cite{capi}. Google Benchmark provides repetition, statistics,
and manual timing for benchmark drivers~\cite{googlebenchmark}.

The prospective contribution is the checked transition from an example's semantic
operation boundary to separate generated execution modes, together with explicit
repetition/rank accounting, original-source navigation, and reproducible experiment
artifacts. The evaluation must establish practical benefits over direct annotations
and conventional drivers. A broader literature review, including TAU, LIKWID,
source-to-source instrumentation, and executable documentation, remains required
before claiming originality or choosing a publication venue.

\section{Limitations and next steps}
The current implementation is a local prototype whose supported semantics should be
made explicit before broader deployment. It assumes repeatable owning-return workloads,
world-collective MPI execution, and CPU completion at function return. It checks the
last result, cannot restore arbitrary mutable state, and records traces whose size
grows with repetitions and annotation density. The Neovim dashboard is one consumer;
the build generator and evaluation script execute without the editor. Remote execution,
accelerator completion rules, bounded-memory aggregation, subcommunicators, and a
state-reset annotation are future extensions rather than completed capabilities.

\bibliographystyle{plain}
\bibliography{references}
\end{document}
